\chapter{Femoral Hernia Repair}
\index{Femoral Hernia Repair}

\section*{Definition and Overview}
A femoral hernia is a bulge of abdominal contents, usually fatty tissue or a loop of bowel, that pushes through a narrow gap in the lower abdominal wall called the femoral canal, located in the groin just below the inguinal ligament. Femoral hernias are much less common than inguinal hernias, but they are more dangerous because the opening is narrow and tight, so the contents are much more likely to become trapped (incarcerated) or have their blood supply cut off (strangulated), which is a surgical emergency. For this reason, femoral hernias are almost always repaired surgically once found, and the operation is called femoral hernia repair. This chapter explains the condition, why prompt repair matters, and what the operation involves.

\section*{Epidemiology}
Femoral hernias are uncommon, accounting for only around 5\% of groin hernias in adults. They occur much more often in women than in men, particularly in older women and in women who have had several pregnancies, though they can occur at any age. Because of their narrow neck, they present with complications more often than other hernias: roughly a third to a half of femoral hernias present as emergencies with incarceration or strangulation. Femoral hernias are rare in children. The risk of strangulation, which can be fatal, is the reason surgical advice is recommended promptly when a femoral hernia is suspected.

\section*{Aetiology and Pathophysiology}
\begin{itemize}
    \item \textbf{Anatomy:} the femoral canal is a small potential space under the inguinal ligament, through which the femoral artery and vein pass to the leg
    \item \textbf{Weakness of the canal:} increased pressure in the abdomen, such as from straining, coughing, heavy lifting, pregnancy, or obesity, pushes contents through the canal
    \item \textbf{Contents:} the hernia sac usually contains fatty tissue (omentum) and sometimes a loop of bowel
    \item \textbf{Incarceration:} the tight opening can trap the contents so they cannot be pushed back, causing pain and a tender lump
    \item \textbf{Strangulation:} if trapped contents include bowel, its blood supply is compressed, leading to ischaemia and gangrene within hours
    \item \textbf{Why women are more affected:} the wider female pelvis and the effects of pregnancy on the abdominal wall predispose women to femoral hernias
    \item \textbf{Emergency risk:} strangulated bowel perforates and causes peritonitis and sepsis, which is life-threatening
\end{itemize}

\section*{Clinical Features}
\begin{itemize}
    \item A swelling or bulge in the upper inner thigh or groin, just below the crease of the groin
    \item The bulge may be small and easy to miss, especially in overweight women
    \item Discomfort or aching in the groin, worse with standing, coughing, or straining
    \item The lump may be reducible (able to be pushed back) early on
    \item Symptoms of incarceration: a painful, tense, non-reducible lump
    \item Symptoms of strangulation: severe pain, nausea, vomiting, abdominal distension, and signs of bowel obstruction
    \item Features of sepsis with strangulation: fever, rapid heart rate, and collapse
\end{itemize}

\section*{Differential Diagnosis}
A groin lump has several possible causes.
\begin{itemize}
    \item Inguinal hernia, which sits higher in the groin and is more common
    \item Enlarged lymph node, which may be tender and can accompany infection
    \item A cyst, including a sebaceous cyst or a swelling of the round ligament
    \item A saphenous varix, a bulging vein at the top of the thigh
    \item A femoral artery aneurysm or pseudoaneurysm
    \item Undescended or retractile testis in the newborn (for inguinal swellings)
    \item Testicular or scrotal conditions such as hydrocele, which extend down to the scrotum
\end{itemize}

\section*{When to Seek Medical Care}
Anyone with a new groin or upper-thigh bulge, especially a woman, should see a doctor promptly. Because femoral hernias so often require urgent surgery, early review is important even if the lump is painless.

\textbf{Contact your local emergency services for:}
\begin{itemize}
    \item A painful, tense groin lump that will not push back
    \item Severe groin or abdominal pain with nausea and vomiting
    \item Fever, sweating, rapid heart rate, or collapse with a groin lump
    \item Abdominal distension and inability to pass wind or stool with a groin lump
\end{itemize}

\section*{Diagnosis and Investigations}
\begin{itemize}
    \item \textbf{History and examination:} a doctor identifies the bulge below the inguinal ligament, and assesses whether it is reducible and tender
    \item \textbf{Ultrasound:} confirms the diagnosis, distinguishes a hernia from other lumps, and assesses the contents
    \item \textbf{CT scan:} used in emergencies or when the diagnosis is uncertain, particularly in strangulation
    \item \textbf{Blood tests:} in emergency presentations, to assess for infection, bowel obstruction, and organ dysfunction
    \item \textbf{Urgent referral:} because of the high risk of strangulation, surgery is planned soon after diagnosis rather than waiting
\end{itemize}

\section*{Management}
\begin{itemize}
    \item \textbf{Elective repair:} once diagnosed, femoral hernias are usually repaired promptly because of the risk of complications
    \item \textbf{Open repair:} an incision in the groin, the hernia contents returned to the abdomen, and the defect closed, often with a mesh patch to strengthen the wall
    \item \textbf{Laparoscopic repair:} keyhole surgery through small incisions, with similar outcomes and a quicker recovery in suitable patients
    \item \textbf{Anaesthesia:} general, spinal, or occasionally local anaesthesia, depending on the patient and approach
    \item \textbf{Emergency surgery:} for incarceration or strangulation, the operation is done urgently; if bowel is dead, the affected segment is removed and the bowel rejoined
    \item \textbf{Recovery:} most patients go home the same day or after a short stay; heavy lifting is avoided for several weeks
    \item \textbf{Pain relief and wound care:} simple analgesia and keeping the wound clean aid recovery
\end{itemize}

\section*{Complications}
\begin{itemize}
    \item Strangulation of bowel with gangrene, perforation, peritonitis, and sepsis if surgery is delayed
    \item Death from delayed strangulation, which is why urgent care matters
    \item Surgical complications: bleeding, infection, wound bruising, and damage to nearby nerves or the femoral vessels
    \item Recurrence of the hernia, which is uncommon with mesh repair
    \item Chronic groin pain in a minority of patients after surgery
    \item Anaesthetic complications, which are rare in healthy patients
\end{itemize}

\section*{Prognosis and Outcomes}
The outlook for femoral hernia repair is excellent when the operation is performed electively, with a low recurrence rate and a quick return to normal activities. When the hernia presents as an emergency, outcomes depend on how quickly surgery occurs and whether the bowel is viable; with prompt treatment, most patients still do well, but the risk of serious complications and death rises substantially with delay. Because femoral hernias are so prone to strangulation, the single most important factor in a good outcome is early diagnosis and repair before complications develop.

\section*{Prevention and Self-Care}
\begin{itemize}
    \item Report any new groin or upper-thigh bulge to a doctor promptly, even if it is painless
    \item Do not attempt to push back a painful, tense lump; seek emergency care
    \item Maintain a healthy weight and avoid chronic straining, which increases abdominal pressure
    \item Treat constipation and chronic cough, which raise the risk of hernia formation and complications
    \item Follow post-operative advice about wound care, lifting restrictions, and gradual return to activity
    \item See a doctor promptly for any recurrence of swelling or pain after repair
\end{itemize}

\section*{Sources and Further Reading}
The following reputable sources provide further information for healthcare students and patients seeking reliable, current advice about femoral hernia repair.
\begin{itemize}
    \item Healthdirect Australia. \textit{Femoral hernia.} https://www.healthdirect.gov.au/hernia
    \item The Royal Australasian College of Surgeons. \textit{Hernia repair patient information.} https://www.surgeons.org/
    \item NHS. \textit{Femoral hernia repair.} https://www.nhs.uk/conditions/femoral-hernia-repair/
    \item Mayo Clinic. \textit{Hernia repair.} https://www.mayoclinic.org/
    \item MSD Manual. \textit{Hernias of the abdominal wall.} https://www.msdmanuals.com/
    \item MedlinePlus. \textit{Femoral hernia.} https://medlineplus.gov/
\end{itemize}
